\documentclass[11pt]{article}
\usepackage[margin=1in]{geometry}
\usepackage{graphicx}
\usepackage{booktabs}
\usepackage{array}
\usepackage{longtable}
\usepackage{amsmath,amssymb}
\usepackage{hyperref}
\usepackage{enumitem}
\usepackage{caption}
\usepackage{float}
\usepackage{verbatim}
\hypersetup{colorlinks=true,linkcolor=blue,citecolor=blue,urlcolor=blue}
\setlength{\parindent}{0pt}
\setlength{\parskip}{0.6em}
\newcommand{\code}[1]{\texttt{#1}}
\newcommand{\model}[1]{\texttt{#1}}
\begin{document}

\begin{center}
{\Large \textbf{A Conservative Benchmark of Explicit Contradiction Detection and Its Limits\\in LLM Hallucination Evaluation Across GPT Model Generations}}\\[1.2em]
Seungho Lee\\
Liberal Arts and Science Academy, 1012 Arthur Stiles Road, Austin, TX 78721, USA\\[0.5em]
Corresponding author: Seungho Lee (seungholee0000@gmail.com); Tel: (512) 414-5272; Fax: (512) 414-6050
\end{center}

\begin{abstract}
Hallucinations--fluent outputs that contain incorrect or unsupported factual claims--remain a primary obstacle to deploying large language models (LLMs) in high-stakes settings. This study evaluates the scope and limits of explicit contradiction detection in HalluDetector, a reference-based detector that compares a model answer to a known correct answer and flags an output only when there is positive evidence of inconsistency, such as contradiction cues, numeric conflicts, or NLI contradiction. This is not a general hallucination or semantic adequacy detector. Across a 150-prompt legacy short-answer benchmark and a 100-prompt enhanced non-atomic benchmark, the updated corpus contains 250 prompt-reference entries and 5,000 archived model answers. For comparability with the legacy run tables, detector-labeled contradiction rates on the legacy short-answer benchmark decrease from 23.5\% for gpt-3.5-turbo-0125 to 14.3\% for gpt-4-0613, 10.3\% for gpt-4o-2024-08-06, and 4.8\% for gpt-5-2025-08-07. The study further evaluates detector behavior on an expanded 850-row, duplicate-aware audit set with two independent annotator labels and adjudicated final labels. On the legacy short-answer audit subset, the detector obtains 84.6\% precision, 94.8\% recall, and 89.4\% F1 for explicit contradiction. On the enhanced non-atomic subset, recall remains high at 90.9\%, but precision falls to 13.5\% and F1 to 23.5\%. This precision drop is treated as a central finding rather than a secondary caveat. A false-positive analysis shows that enhanced errors are concentrated in uncertainty-oriented, source-limited, and cautionary prompts, where safe answers often mention unavailable evidence, claims to avoid, or limits on what can be concluded. Overall, HalluDetector is best interpreted as a recall-oriented reference-based contradiction screen: it performs best on compact answer-key settings and exposes a clear limitation when calibrated uncertainty or source-limitation language resembles contradiction evidence in non-atomic prompts.
\end{abstract}

\textbf{Keywords:} contradiction detection; explicit contradiction; large language models; factuality; hallucination evaluation; natural language inference; benchmark evaluation; model auditing; abstention; reproducibility; AI safety

\section*{Introduction}
Hallucinations in neural language generation are commonly defined as outputs that are not supported by source grounding, reference truth, or the model's training distribution. In modern LLM assistants, hallucinations are particularly risky because they are often delivered with a confident tone and detailed justification, encouraging over-trust. This has motivated benchmarks and interventions targeting truthfulness and misconception resistance. TruthfulQA, for instance, demonstrated that models can systematically mimic human-written misinformation and misconceptions, and that raw scaling can amplify these effects if alignment is not prioritized (3). Instruct tuning and reinforcement learning from human feedback (RLHF) have been shown to improve instruction-following and reduce misleading behavior (2). The GPT-4 technical report similarly documents improvements in factuality and reduced hallucination relative to GPT-3.5 in internal evaluations (1).

This paper studies contradiction-style hallucination behavior in a controlled, reference-based setting, where every prompt has a known correct reference answer. This enables stable measurement, interpretability, and debugging, but it also exposes a central limitation: reference-based contradiction detection can mistake cautious discussion of uncertainty or source limits for contradiction evidence. The accompanying repository, HalluDetector, provides: (i) prompt sets spanning factual, riddle, auto-generated adversarial, and enhanced non-atomic questions; (ii) an explicit contradiction detector with documented decision rules and thresholds; and (iii) end-to-end scripts to generate model answers, label them, compute metrics, and render graphs. The revised analysis therefore evaluates both where the detector works and where its scope limits appear.

\textbf{Scope and definition.} This study focuses on reference-based hallucination detection: given a model answer $a$ and a reference answer $a^*$, the detector labels $a$ as hallucinated if it contains positive evidence of falsity relative to $a^*$. A key design choice of HalluDetector is that incompleteness is not hallucination. If $a$ fails to include the correct content but also does not contradict it, the detector does not label hallucination. This policy can under-count semantically inadequate answers that remain non-contradictory with respect to the reference. The revised audit therefore records explicit contradiction separately from \code{incomplete\_or\_vacuous} responses, abstentions, unsupported extra claims, and needs-review cases, allowing the analysis to distinguish contradiction detection from broader semantic adequacy evaluation.

\textbf{Contributions.} This paper makes five contributions: (1) it formalizes the detector and thresholds in mathematical form; (2) it introduces a conservative, fully auditable contradiction-evidence benchmark pipeline that separates explicit falsity from mere incompleteness; (3) it reports comparative detector-labeled explicit contradiction rates across four GPT endpoint identifiers and five runs each; (4) it evaluates detector performance against a two-annotator adjudicated audit set while separately reporting explicit contradiction labels, semantic-failure labels, duplicate-aware statistics, and prompt-level aggregation; and (5) it analyzes the enhanced non-atomic precision drop through prompt-category breakdowns, a false-positive taxonomy, and component-level attribution.

\section*{Related Work}
\textbf{Truthfulness and hallucination benchmarks.} TruthfulQA introduced adversarial questions built around common misconceptions, showing that many LLMs produce fluent falsehoods and that model size can amplify imitation of misinformation (3). These findings motivated additional evaluation suites and alignment-focused approaches to improve truthfulness. More recent empirical work further studies factuality hallucination in large language models and expands the benchmark landscape (9).

\textbf{Alignment and post-training.} InstructGPT and related RLHF systems demonstrated that human feedback can steer models toward better instruction-following, reduced toxic content, and more truthful behavior (2). OpenAI's GPT-4 technical report documents improvements in factuality and reduced hallucination on internal evaluations (1). While this study does not attempt to attribute causality to specific training methods, it reports lower detector-labeled contradiction-style hallucination rates across the evaluated endpoints.

\textbf{Reference-based textual factuality.} In summarization, FactCC and follow-on work framed factual consistency as a classification problem using synthetic perturbations and NLI signals (5). In claim verification, FEVER established a benchmark for verifying claims against evidence, catalyzing retrieval + NLI pipelines (4). The methodological core--entailment vs. contradiction between generated claims and trusted evidence--is shared with the approach used here. This study differs in packaging these ideas into a compact, reproducible benchmark with explicit detector thresholds, cross-run endpoint comparisons, and a conservative abstention-aware labeling policy. Recent surveys also synthesize broader hallucination and factuality evaluation trends in large language models (10, 11).

\textbf{Zero-resource detectors.} When references are unavailable, black-box detectors such as SelfCheckGPT measure hallucination risk through generation variance and cross-sample inconsistency (6). These methods can be deployed without external evidence but increase inference cost and can fail when a model confidently and consistently repeats the same falsehood.

\textbf{Semantic similarity and embeddings.} Embedding similarity is widely used to evaluate paraphrases and to support flexible matching beyond exact string equality. Sentence-BERT and Sentence Transformers provide practical sentence-level embeddings for such comparisons (7). HalluDetector uses embedding cosine similarity as an auxiliary baseline signal with guardrails to reduce short-text false positives.

\textbf{Reference selection.} The reference list combines foundational sources on truthfulness, claim verification, NLI, and embedding-based semantic similarity with recent 2024--2025 work on factuality and hallucination evaluation. Citations are included only where they directly support the surrounding methodological or interpretive claim.

\section*{Materials and Methods}
\textbf{Problem setup.} Given a prompt $q$, reference answer $a^*$, and model response $a$, the detector outputs a binary label $H(a,a^*) \in \{0,1\}$ indicating whether $a$ contains hallucinated content relative to $a^*$. HalluDetector adopts a contradiction-evidence policy:

\begin{quote}
Hallucination requires evidence of inconsistency. Paraphrases are allowed; incompleteness is not labeled hallucination unless the response introduces a contradiction or a mutually exclusive alternative.
\end{quote}

This policy is formalized as a disjunction of contradiction evidence signals:
\begin{equation}
H(a,a^*) = \mathbf{1}[C_{cue} \lor C_{nli} \lor C_{num} \lor C_{alt} \lor C_{unavail} \lor C_{atomic}].
\tag{eq1}
\end{equation}

\textbf{Normalization.} Both $a$ and $a^*$ are canonicalized via Unicode normalization, lowercasing, whitespace normalization, and punctuation stripping that preserves scientific characters (e.g., - . / \^{} + : =). This reduces formatting-induced mismatches while keeping numerically meaningful tokens.

\textbf{Lexical shortcut signals.} The detector computes three lexical agreement signals:
\begin{itemize}[leftmargin=1.5em]
\item Exact match: $C_{exact}=\mathbf{1}[canon(a)=canon(a^*)]$.
\item Substring match: $C_{substr}=\mathbf{1}[canon(a^*)\subseteq canon(a)]$.
\item Token subsequence match: a weaker containment check used as a non-veto support signal.
\end{itemize}
These checks are not sufficient to declare correctness in open-world settings; in this reference-based setting they serve as strong evidence of agreement, but other contradiction checks can still veto.

\textbf{Numeric consistency and conflict-if-present.} Numeric hallucinations are common in knowledge Q\&A (e.g., years, dataset IDs). If the reference is numeric or coordinate-like, the detector parses numeric values and defines a tolerant match by
\begin{equation}
|x-x^*| \leq \max(\delta_{abs}, \delta_{rel}|x^*|),
\tag{eq2}
\end{equation}
with defaults $\delta_{rel}=5\times 10^{-3}$ and $\delta_{abs}=5\times 10^{-2}$. For integer-like values with magnitude $\geq 100$ (years/IDs), exact match is required. If $a$ contains numbers but none matches the reference target, the detector treats this as contradiction evidence:
\begin{equation}
C_{num}=\mathbf{1}[numbers(a)\neq\emptyset \land \forall x\in numbers(a), x\not\approx x^*].
\tag{eq3}
\end{equation}
If $a$ contains no numbers, the detector does not label hallucination, consistent with the incompleteness policy. If $a$ provides multiple mutually exclusive numeric alternatives, the detector flags hallucination:
\begin{equation}
C_{alt}=\mathbf{1}[alt\_marker(a) \land multi\_number(a)].
\tag{eq4}
\end{equation}

\textbf{Embedding similarity baseline.} To support paraphrases, HalluDetector computes cosine similarity between sentence embeddings:
\begin{equation}
s_{cos}(a,a^*) = \frac{\langle e(a),e(a^*)\rangle}{\|e(a)\|\|e(a^*)\|},
\tag{eq5}
\end{equation}
with encoder $e(\cdot)=\code{sentence-transformers/all-MiniLM-L6-v2}$ and threshold $\tau_{emb}=0.72$. Guardrails require at least a minimal overlap of content tokens for very short gold answers. If embeddings are unavailable, a lexical Jaccard fallback is used:
\begin{equation}
J(a,a^*)=\frac{|T(a)\cap T(a^*)|}{|T(a)\cup T(a^*)|}, \qquad J(a,a^*)\geq \tau_J, \quad \tau_J=0.40.
\tag{eq6}
\end{equation}

\textbf{Bidirectional NLI contradiction veto.} The detector uses an MNLI-trained NLI model (default \code{roberta-large-mnli}) (8) to compute entailment and contradiction probabilities in both directions:
\begin{equation}
p^{\rightarrow}_{contr}=P(contr\mid a\Rightarrow a^*), \qquad p^{\leftarrow}_{contr}=P(contr\mid a^*\Rightarrow a),
\tag{eq7}
\end{equation}
and likewise for entailment probabilities $p^{\rightarrow}_{ent}$ and $p^{\leftarrow}_{ent}$. A contradiction veto triggers when
\begin{equation}
C_{nli}=\mathbf{1}[\max(p^{\rightarrow}_{contr},p^{\leftarrow}_{contr})\geq \tau_c \land \min(p^{\rightarrow}_{ent},p^{\leftarrow}_{ent})\leq \tau_e],
\tag{eq8}
\end{equation}
with defaults $\tau_c=0.75$ and $\tau_e=0.30$.

\textbf{Unavailability and atomic mismatch heuristics.} If the model response asserts non-existence or unavailability while the gold is concrete, the detector flags hallucination unless other agreement signals hold:
\begin{equation}
C_{unavail}=\mathbf{1}[unavailability\_cue(a) \land \neg C_{substr} \land \neg C_{exact} \land \neg numeric\_match].
\tag{eq9}
\end{equation}
For atomic gold answers, if the response is also short but shares no content tokens and does not match numerically, the detector flags hallucination:
\begin{equation}
C_{atomic}=\mathbf{1}[atomic(a^*) \land short(a) \land no\_shared\_content(a,a^*)].
\tag{eq10}
\end{equation}

\textbf{Detector pipeline diagram.} Figure~\ref{fig:detector-pipeline} shows the detector pipeline.
\begin{figure}[H]
\centering
\includegraphics[width=0.62\linewidth]{figures/detector_pipeline_flowchart_diagram.png}
\caption{Flowchart of the HalluDetector detector pipeline.}
\label{fig:detector-pipeline}
\end{figure}

\section*{Experimental Setup}
\textbf{Prompt sets.} The legacy short-answer benchmark is stored as three prompt CSVs (per run) with schema \code{template, id, prompt, correct\_answer}: factual (50), riddle (50), and auto-generated (50). The updated corpus additionally includes a 100-prompt enhanced non-atomic benchmark covering longer-form explanation, causal reasoning, synthesis, citation-grounded answering, and uncertainty-oriented prompts. The enhanced prompt set is organized evenly across five templates (20 prompts per template category): long-form explanation, causal reasoning, synthesis, citation-grounded answering, and nuanced uncertainty. Across both phases, the prompt inventory contains 250 prompt-reference entries. The audit uses benchmark phase, source category/file, and prompt ID as the underlying prompt identity so that run-specific wording variation in auto-generated prompts does not inflate prompt-level counts. The enhanced prompts therefore form a substantial stress-test expansion of the legacy short-answer benchmark, while the manuscript continues to treat the corpus as compact rather than as a comprehensive broad hallucination benchmark.


\begin{table}[H]
\centering
\scriptsize
\caption*{Representative prompt examples from the legacy short-answer and enhanced non-atomic prompt sets. The legacy short-answer set contains compact answers, but many prompts require exact reasoning, ambiguity resolution, symbolic manipulation, or calculation rather than simple factual lookup. The complete prompt database is provided in the DOI-linked materials.}
\begin{tabular}{p{0.20\linewidth}p{0.53\linewidth}p{0.20\linewidth}}
\toprule
Prompt group & Representative prompt & Main test\\
\midrule
Legacy logic/calendar & Across a full 400-year Gregorian cycle, how many months have five Saturdays and five Sundays? & Calendar arithmetic and exact counting\\
Legacy physics/logic & Feathers and gold have equal true mass in vacuum; which reads heavier on a spring scale in air and why? & Buoyancy and apparent weight\\
Legacy math & How many 5-card poker hands contain exactly two face cards and no aces? & Combinatorics\\
Legacy language & Combine MINUTE and MOMENT; which letters occur an odd number of times? & Symbolic word reasoning\\
Legacy riddle/logic & What gets larger when you take away part of it? & Nonliteral reasoning\\
\addlinespace[0.15em]
Enhanced long-form explanation & Explain why increasing sample size reduces random error but does not automatically remove bias. & Multi-sentence explanation\\
Enhanced causal reasoning & A new tutoring program begins and student test scores rise. Why is this not enough to prove the tutoring caused the improvement? & Causation vs. correlation\\
Enhanced synthesis & Compare bacteria and viruses in relation to antibiotic treatment. & Combining source-grounded claims\\
Enhanced citation-grounded answering & Using only this source: The trial improved symptoms but did not include a placebo control group. What limitation should be mentioned? & Source-limited answering\\
Enhanced nuanced uncertainty & A user asks for the current population of a city but your only source is ten years old. What should a reliable answer say? & Cautious uncertainty\\
\bottomrule
\end{tabular}
\end{table}

\textbf{Model endpoints and run protocol.} The study evaluates four endpoints as referenced by output directory conventions: gpt-3.5-turbo-0125, gpt-4-0613, gpt-4o-2024-08-06, and gpt-5-2025-08-07. Each endpoint is evaluated five times on the legacy 150-prompt benchmark and five times on the enhanced 100-prompt benchmark, producing 3,000 legacy-benchmark answers and 2,000 enhanced-benchmark answers. The legacy-benchmark runs were conducted on March 11 between 3:00 PM and 4:00 PM EST. Enhanced-run timestamps are preserved in the archived enhanced output directories. Model answers are generated through the OpenAI API via a response-generation script, cleaned to a concise first-line form, and labeled by the detector with stable thresholds.

\textbf{Decoding defaults.} The generation defaults are known from the archived generation helper rather than inferred after the fact. For HuggingFace causal-language-model runs, the helper defaults to \code{max\_new\_tokens=50} with greedy decoding unless explicit sampling arguments are supplied. For OpenAI runs, non-GPT-5 endpoints use Chat Completions and preserve legacy family defaults: gpt-3.5-turbo-0125 uses temperature 1.0 and max tokens 512; gpt-4-0613 uses temperature 0.8 and max tokens 2000; and gpt-4o-2024-08-06 uses temperature 0.8 and max tokens 4000. The gpt-5-2025-08-07 endpoint is generated through the Responses API with \code{max\_output\_tokens=512}; the helper intentionally does not pass temperature or top-p to this GPT-5 variant unless such arguments are explicitly supported by the model family. All results in this paper are those recorded in the repository outputs.

\textbf{Evaluation metrics.} The analysis computes detector-labeled contradiction rate by model, category, and run; run-to-run variability; detector precision/recall/F1 against audit labels; and prompt-level pattern summaries via repository graphs. For uncertainty, it reports 95\% prompt-cluster bootstrap confidence intervals for overall detector-labeled contradiction rates. Each bootstrap replicate resamples prompt IDs with replacement within each category, then aggregates all five runs for the sampled IDs.

For the revised audit analysis, the manuscript additionally reports duplicate-aware statistics. The audit set is evaluated at three levels: (i) row-level metrics over all 850 audited rows; (ii) deduplicated prompt-response metrics over 761 unique prompt-response pairs; and (iii) prompt-identity aggregation over 242 underlying prompt identities represented in the audit. Audit labels use a multi-class codebook: \code{supported}, \code{contradiction}, \code{incomplete\_or\_vacuous}, \code{abstention}, \code{unsupported\_extra\_claim}, and \code{needs\_review}. Precision, recall, and F1 are binary explicit-contradiction metrics: \code{contradiction} is the positive class, and all other audit labels are negative and reported separately.

\section*{Results and Discussion}
\textbf{Detector-labeled contradiction rates by model, category, and run.} Table 1 reports detector-labeled contradiction rates for the legacy 150-prompt benchmark for each model in each run. The results show lower detector-labeled explicit contradiction rates from gpt-3.5-turbo-0125 to gpt-4-0613 to gpt-4o-2024-08-06, with gpt-5-2025-08-07 having the lowest rate on this legacy benchmark. Because the study evaluates recorded API outputs associated with those endpoint identifiers, this is interpreted as a benchmark-specific comparison of behavior during the collection period. The remaining detector-positive cases concentrate in numeric/database-style questions and adversarial auto-generated prompts.

\begin{table}[H]
\centering
\small
\caption{Detector-labeled contradiction rates (\%) by model, category, and run on the legacy short-answer benchmark. Each run contains 150 prompts (50 per category).}
\begin{tabular}{llrrrr}
\toprule
Model & Run & Factual & Riddle & Auto-gen & Overall\\
\midrule
gpt-3.5-turbo-0125 & gpt3.5-turbo test & 18.0\% & 30.0\% & 22.0\% & 23.3\%\\
 & gpt3.5-turbo test2 & 20.0\% & 28.0\% & 24.0\% & 24.0\%\\
 & gpt3.5-turbo test3 & 24.0\% & 30.0\% & 20.0\% & 24.7\%\\
 & gpt3.5-turbo test4 & 22.0\% & 30.0\% & 14.0\% & 22.0\%\\
 & gpt3.5-turbo test5 & 14.0\% & 28.0\% & 28.0\% & 23.3\%\\
gpt-4-0613 & gpt4 test & 6.0\% & 24.0\% & 8.0\% & 12.7\%\\
 & gpt4 test2 & 8.0\% & 26.0\% & 20.0\% & 18.0\%\\
 & gpt4 test3 & 8.0\% & 26.0\% & 2.0\% & 12.0\%\\
 & gpt4 test4 & 6.0\% & 26.0\% & 10.0\% & 14.0\%\\
 & gpt4 test5 & 8.0\% & 26.0\% & 10.0\% & 14.7\%\\
gpt-4o-2024-08-06 & gpt4o test & 8.0\% & 16.0\% & 0.0\% & 8.0\%\\
 & gpt4o test2 & 6.0\% & 16.0\% & 2.0\% & 8.0\%\\
 & gpt4o test3 & 10.0\% & 20.0\% & 8.0\% & 12.7\%\\
 & gpt4o test4 & 8.0\% & 18.0\% & 4.0\% & 10.0\%\\
 & gpt4o test5 & 12.0\% & 18.0\% & 8.0\% & 12.7\%\\
gpt-5-2025-08-07 & gpt-5-2025-08-07 test & 10.0\% & 0.0\% & 8.0\% & 6.0\%\\
 & gpt-5-2025-08-07 test2 & 6.0\% & 0.0\% & 0.0\% & 2.0\%\\
 & gpt-5-2025-08-07 test3 & 8.0\% & 0.0\% & 4.0\% & 4.0\%\\
 & gpt-5-2025-08-07 test4 & 12.0\% & 2.0\% & 8.0\% & 7.3\%\\
 & gpt-5-2025-08-07 test5 & 8.0\% & 0.0\% & 6.0\% & 4.7\%\\
\bottomrule
\end{tabular}
\end{table}

\textbf{Inter-run variability.} Most variability arises on the auto-generated prompt set. This is consistent with two plausible mechanisms: prompts closer to the decision boundary and sensitivity to early token choices after first-line cleaning. Notably, the factual and riddle categories remain comparatively stable, while most variability is concentrated in the auto-generated category.

\textbf{Abstention as a formal result.} Empty first-line answers are counted by the generation script as non-hallucinated abstentions rather than contradictions. Abstention is concentrated in gpt-5-2025-08-07: across all 750 archived answers for that endpoint identifier, 146 are empty after first-line cleaning (19.5\%), versus 0 for the other evaluated endpoints. The abstention rate is especially high on riddles (30.8\% averaged across runs), indicating that part of the reduction in detector-labeled contradiction-style hallucination reflects more conservative answering behavior rather than only improved task success.

\textbf{Detector component ablation.} To make the detector decision path more explicit, the archived-output ablation keeps only contradiction cues and numeric conflict checks in one stage, adds unavailability and atomic-mismatch heuristics in another, and compares these with the archived detector output. Embedding similarity is not shown as a separate decision-stage ablation because in this implementation it supports agreement decisions but does not independently force hallucination labels.

\subsection*{Detector performance vs. adjudicated audit}
The adjudicated audit contains 850 rows: 450 rows from the legacy short-answer audit and 400 rows from the enhanced audit sample. The enhanced sample therefore represents 47.1\% of the row-level audit and was added specifically to test prompts less favorable to short-answer contradiction scoring. Each row contains the prompt, reference answer, model response, detector label, two annotator labels, annotator rationales, agreement status, adjudicated label, duplicate identifiers, and prompt-cluster identifiers. Annotators were blinded to detector output and model identity. Disagreements were resolved by adjudication using the written label codebook.

The audit label set separates explicit contradiction from other forms of poor answer quality. In particular, \code{incomplete\_or\_vacuous} responses are not treated as successful answers, even when they do not explicitly contradict the reference. Likewise, \code{unsupported\_extra\_claim} is reserved for responses that provide a correct core answer but add an inaccurate or unsupported additional claim. This label structure directly addresses the limitation that a contradiction-only detector can mark vacuous or incomplete outputs as non-hallucinations. Borderline cases are adjudicated according to the explicit-contradiction codebook: reasonable paraphrases, correct core answers, and underspecified but non-mutually-exclusive answers are treated as non-contradictions, while direct conflicts with the reference remain labeled as contradictions.

Across the 850 adjudicated audit rows, 717 responses are labeled \code{supported}, 69 \code{contradiction}, 57 \code{incomplete\_or\_vacuous}, 3 \code{abstention}, 3 \code{unsupported\_extra\_claim}, and 1 \code{needs\_review}. The two annotators agree on 811 of 850 rows (95.4\%), with Cohen's $\kappa = 0.845$. The 39 disagreements primarily involve the boundary between \code{supported} and \code{incomplete\_or\_vacuous} for short or partially adequate responses. The adjudicated labels are used as the final audit labels in Table 2.

Table 2 reports detector performance with explicit contradiction as the positive class. The legacy short-answer audit subset is the closest apples-to-apples detector-quality comparison with the legacy benchmark. The enhanced audit subset is a stress test of scope: because it includes longer-form, causal, synthesis, citation-grounded, and uncertainty-oriented prompts, the detector more often over-flags short but supportable paraphrases, cautious non-answers, or underspecified answers as contradiction-like. The resulting 13.5\% precision and 23.5\% F1 on the enhanced subset are therefore not treated as minor secondary results. They identify a boundary condition for the method: HalluDetector remains recall-oriented for explicit contradictions, but it is not sufficient by itself for scoring non-atomic prompts where calibrated answers discuss uncertainty, evidence limits, or claims that should be avoided. For this reason, the combined overall metric is reported for transparency but is not interpreted as a replacement for either the legacy short-answer result or the enhanced stress-test analysis.

A manual review of the binary error rows supports this interpretation. Of the 74 row-level false positives, 64 come from the enhanced non-atomic audit subset, and 49 come from uncertainty-oriented prompts. These prompts often ask what a reliable answer should avoid or what should be said when evidence is unavailable, incomplete, or outdated. Correct or cautious responses can therefore mention phrases such as unavailable information, unverified claims, or claims to avoid. The detector sometimes treats those cautionary mentions as contradiction evidence even when the adjudicated answer is supported or abstaining rather than contradictory. This explains why enhanced-subset precision is low without implying that the legacy short-answer benchmark result is invalid. A direct detector-level follow-up is to add an uncertainty-aware guardrail that distinguishes explicit factual contradiction from cautious statements about unavailable, unsupported, or insufficient evidence.

\begin{table}[H]
\centering
\small
\caption*{Binary interpretation of adjudicated audit labels for the explicit-contradiction metric. The negative labels are still reported because they identify semantic or completeness problems outside the detector's binary target.}
\begin{tabular}{p{0.24\linewidth}p{0.25\linewidth}p{0.43\linewidth}}
\toprule
Adjudicated label & Binary metric role & Interpretation\\
\midrule
\code{contradiction} & Positive & Direct conflict with the reference or source constraint.\\
\code{supported} & Negative & Correct, paraphrased, or underspecified but not mutually inconsistent.\\
\code{incomplete\_or\_vacuous} & Negative, separately reported & Empty, vague, or insufficient response without explicit contradiction.\\
\code{abstention} & Negative, separately reported & Refusal or non-answer rather than a conflicting factual claim.\\
\code{unsupported\_extra\_claim} & Negative, separately reported & Correct core answer plus unsupported extra content not treated as explicit contradiction.\\
\code{needs\_review} & Negative, separately reported & Ambiguous row retained for transparency.\\
\bottomrule
\end{tabular}
\end{table}

\begin{table}[H]
\centering
\small
\caption{Detector performance on the 850-row adjudicated audit. TP/FP/FN are binary explicit-contradiction metrics: adjudicated \code{contradiction} is positive and all other adjudicated labels are negative.}
\begin{tabular}{lrrrrrrrr}
\toprule
Audit subset & $n$ & TP & FP & FN & TN & Precision & Recall & F1\\
\midrule
Legacy short-answer audit & 450 & 55 & 10 & 3 & 382 & 84.6\% & 94.8\% & 89.4\%\\
Enhanced non-atomic audit & 400 & 10 & 64 & 1 & 325 & 13.5\% & 90.9\% & 23.5\%\\
Overall row-level audit & 850 & 65 & 74 & 4 & 707 & 46.8\% & 94.2\% & 62.5\%\\
\bottomrule
\end{tabular}
\end{table}

To account for repeated prompts and repeated prompt-response pairs, Table 3 reports duplicate-aware audit statistics. The audit contains 242 underlying prompt identities represented in the sampled audit and 761 unique prompt-response pairs. A naive text-hash count is higher because some auto-generated prompt IDs have run-specific wording variants; therefore, prompt-level aggregation uses benchmark phase, source category/file, and prompt ID rather than raw prompt text alone. Row-level metrics are reported alongside deduplicated prompt-response metrics and prompt-identity aggregation. Prompt-cluster bootstrap confidence intervals for the combined row-level explicit-contradiction metrics are 32.6\%--62.3\% for precision, 88.6\%--98.6\% for recall, and 48.4\%--75.1\% for F1.

\begin{table}[H]
\centering
\scriptsize
\caption{Duplicate-aware audit statistics.}
\resizebox{\linewidth}{!}{%
\begin{tabular}{lrrrrrrrrrr}
\toprule
Evaluation level & Rows/items & Prompt identities & Unique prompt-responses & TP & FP & FN & TN & Precision & Recall & F1\\
\midrule
Row-level audit & 850 & 242 & 761 & 65 & 74 & 4 & 707 & 46.8\% & 94.2\% & 62.5\%\\
Deduplicated prompt-response & 761 & 242 & 761 & 65 & 73 & 4 & 619 & 47.1\% & 94.2\% & 62.8\%\\
Prompt-identity aggregation & 242 & 242 & -- & 38 & 30 & 2 & 172 & 55.9\% & 95.0\% & 70.4\%\\
\bottomrule
\end{tabular}}
\end{table}

\subsection*{Enhanced non-atomic stress-test analysis}
The enhanced non-atomic subset is not used as evidence of broad generalization; instead, it exposes the detector's main scope limitation. Table~\ref{tab:enhanced-category-breakdown} separates the enhanced audit by prompt category. Because the enhanced audit sample is stratified by detector and audit behavior rather than category-balanced, row counts differ across the five enhanced prompt types. Within this enhanced audit sample, detector precision and F1 are highest on citation-grounded prompts, where the source constraint often creates a clearer contradiction target. In contrast, uncertainty-oriented prompts produce 49 false positives and only 3 true positives. This yields 5.8\% precision and 10.9\% F1 for that category even though recall is 100.0\%. The main source of low precision is therefore not missed contradictions; it is over-triggering when safe answers discuss evidence limits, missing sources, or claims to avoid. Categories with no adjudicated contradiction positives have undefined recall and F1, so their zero precision should be read as false-positive evidence rather than as a complete performance estimate.

\begin{table}[H]
\centering
\scriptsize
\caption{Prompt-category breakdown within the enhanced non-atomic audit subset. Precision, recall, and F1 are binary explicit-contradiction metrics. Recall and F1 are left undefined when a category has no adjudicated contradiction positives.}
\label{tab:enhanced-category-breakdown}
\begin{tabular}{lrrrrrrrr}
\toprule
Enhanced category & $n$ & TP & FP & FN & TN & Precision & Recall & F1\\
\midrule
Long-form explanation & 103 & 0 & 5 & 0 & 98 & 0.0\% & -- & --\\
Causal reasoning & 69 & 0 & 2 & 0 & 67 & 0.0\% & -- & --\\
Synthesis & 77 & 0 & 4 & 0 & 73 & 0.0\% & -- & --\\
Citation-grounded answering & 38 & 7 & 4 & 1 & 26 & 63.6\% & 87.5\% & 73.7\%\\
Nuanced uncertainty & 113 & 3 & 49 & 0 & 61 & 5.8\% & 100.0\% & 10.9\%\\
\bottomrule
\end{tabular}
\end{table}

The 64 enhanced false positives were then coded into mutually exclusive explanation classes by reviewing the prompt, reference answer, model response, and adjudicated label. Table~\ref{tab:fp-taxonomy} shows that false positives are not randomly distributed. Thirty involve source availability, current-data limits, missing documents, missing charts, future events, or other evidence-limitation language. Nineteen involve a safe answer naming the prohibited action or claim, such as giving a diagnosis, treatment plan, confident identification, or unsupported medical advice. Fifteen are supportable paraphrases or compressed answers in explanation, synthesis, causal, or citation-grounded prompts. This taxonomy supports the central interpretation: contradiction detection can penalize exactly the caution and calibration behavior that broader hallucination evaluation often wants to reward.

\begin{table}[H]
\centering
\small
\caption{False-positive taxonomy for the 64 enhanced non-atomic detector positives that were adjudicated non-contradictory. Categories are mutually exclusive.}
\label{tab:fp-taxonomy}
\begin{tabular}{p{0.34\linewidth}rp{0.48\linewidth}}
\toprule
False-positive class & Count & Mechanism\\
\midrule
Source, availability, or current-evidence limitation & 30 & Safe answers stated that data were unavailable, outdated, not visible, not provided, or required current verification.\\
Forbidden-claim or cautionary-action mention & 19 & Safe answers named the claim or action that should be avoided, such as diagnosis, treatment advice, or confident identification.\\
Supportable paraphrase or compressed answer & 15 & Correct or acceptable paraphrases were treated as contradiction-like because they omitted wording, compressed the reference, or used a different explanation structure.\\
\bottomrule
\end{tabular}
\end{table}

The detector core returns rich per-row diagnostics, including a final reason code, lexical and numeric flags, embedding diagnostics, NLI probabilities, contradiction-veto status, unavailability flags, and weak-signal counts. In the historical response-generation CSVs, only summary columns were written. Table~\ref{tab:component-reconstruction} therefore reports component attribution from the stored rows, the detector decision order, and reproducible rule checks. The deterministic re-run reproduced 18 enhanced false positives as unavailability-claim conflicts with the reference. The remaining 46 detector-positive false positives did not match deterministic numeric, atomic-mismatch, or exact lexical-contradiction rules in the reconstruction, so they are grouped as NLI/semantic contradiction-trigger cases for analysis.

\begin{table}[H]
\centering
\small
\caption{Component attribution for enhanced false positives.}
\label{tab:component-reconstruction}
\begin{tabular}{p{0.38\linewidth}rp{0.43\linewidth}}
\toprule
Detector component or trigger class & Count & Interpretation\\
\midrule
Unavailability / nonexistence heuristic & 18 & Reproduced by deterministic rules; cautious statements such as unavailable or not available were treated as conflicting with a concrete reference.\\
Inferred NLI / semantic contradiction trigger & 46 & Detector-positive rows not explained by deterministic numeric, atomic, or unavailability checks; these include caveats, paraphrases, and source-limitation language.\\
Numeric conflict or numeric alternatives & 0 & Not a driver of enhanced false positives.\\
Atomic entity mismatch & 0 & Not a driver of enhanced false positives.\\
Explicit lexical contradiction cue & 0 & Not a driver of enhanced false positives in the reconstructed pass.\\
\bottomrule
\end{tabular}
\end{table}

These analyses shift the interpretation of the enhanced benchmark. The enhanced subset shows that contradiction-evidence detectors can perform well when the target is a compact reference disagreement but lose precision when prompts ask models to reason about the limits of evidence. In those settings, a reliable answer may explicitly mention that a claim is unsupported, a source is missing, or a conclusion should not be drawn. A detector that only looks for contradiction cues can confuse this metalinguistic caution with a factual contradiction. This pattern is scientifically meaningful: explicit contradiction detection is entangled with uncertainty expression, source limitation, abstention, partial correctness, and explanation structure.

\subsection*{Error analysis (audit disagreements)}
The adjudicated audit changes the interpretation of detector errors. In the legacy short-answer subset, remaining errors continue to involve synonym, notation, and granularity issues, such as equivalent scientific notation or overly broad chemical names. In the enhanced subset, the taxonomy above shows that many false positives arise from cautionary or source-limitation language rather than direct factual conflict. These cases show over-sensitivity in the binary detector rather than invalidating the legacy short-answer result.

The audit therefore distinguishes two failure modes that were previously conflated: (i) explicit contradiction errors, where the model response conflicts with the reference; and (ii) semantic adequacy errors, where the response is empty, vague, incomplete, or insufficiently grounded but not explicitly contradictory. Semantic-adequacy failures remain important, but they are reported through separate labels rather than treated as positives in the binary explicit-contradiction metric. This distinction is important because HalluDetector was designed as a conservative contradiction detector, not a general answer-quality evaluator.

False positives in the adjudicated audit occur when the detector flags contradiction but adjudication assigns any non-contradiction label. In the current audit, these false positives are mostly \code{supported} responses rather than semantic-failure labels: 71 of 74 row-level false positives are adjudicated \code{supported}, and the remaining 3 are \code{abstention}. False negatives occur when adjudication identifies explicit contradiction but the detector does not. The adjudicated audit contains 74 row-level false positives and 4 row-level false negatives overall. The legacy short-answer subset accounts for 10 false positives and 3 false negatives; the enhanced subset accounts for 64 false positives and 1 false negative. This pattern shows that the detector remains recall-oriented for explicit contradictions, while the enhanced subset reveals over-sensitivity to terse but supportable answers and cautious non-answers. Non-contradictory adequacy failures are therefore kept separate rather than folded into the binary explicit-contradiction metric.

\subsection*{Prompt-template patterns and visuals}
The repository includes graphs per run: \code{hallucinations by template.png}, \code{hallucinations by keywords.png}, and \code{feature correlation heatmap.png}. Figures 2--4 show the visual summaries for gpt-3.5-turbo-0125 run gpt3.5-turbo test.

\begin{figure}[H]
\centering
\includegraphics[width=0.85\linewidth]{figures/page009_img001.png}
\caption{Detector-labeled contradiction-style hallucination rate by template for gpt-3.5-turbo-0125 (run gpt3.5-turbo test).}
\end{figure}

\begin{figure}[H]
\centering
\includegraphics[width=0.85\linewidth]{figures/page010_img001.png}
\caption{Detector-labeled contradiction-style hallucination rate by keyword for gpt-3.5-turbo-0125 (run gpt3.5-turbo test).}
\end{figure}

\begin{figure}[H]
\centering
\includegraphics[width=0.70\linewidth]{figures/page010_img002.png}
\caption{Feature correlation heatmap for gpt-3.5-turbo-0125 (run gpt3.5-turbo test). Note: the repo defines ``is correct'' as ``not hallucinated'' for plotting convenience.}
\end{figure}

Table 7 summarizes template-level extremes for a representative early model and for the latest model run.

\begin{table}[H]
\centering
\caption{Top detector-labeled contradiction-style hallucination-rate templates in gpt-3.5-turbo-0125 (run gpt3.5-turbo test) vs gpt-5-2025-08-07 (run gpt-5-2025-08-07 test5). Rates are detector-labeled.}
\label{tab:template-extremes}
\begin{tabular}{lrlr}
\toprule
\multicolumn{2}{c}{gpt-3.5-turbo-0125} & \multicolumn{2}{c}{gpt-5-2025-08-07}\\
\midrule
Template & Rate & Template & Rate\\
\midrule
finance & 100.0\% & chemistry & 25.0\%\\
math & 50.0\% & geography & 14.3\%\\
computer science & 50.0\% & physics & 5.9\%\\
geography & 42.9\% & logic & 5.3\%\\
biology & 33.3\% & astronomy & 0.0\%\\
science & 33.3\% & economics & 0.0\%\\
\bottomrule
\end{tabular}
\end{table}

\textbf{Length correlations.}

Across all runs, prompt length and response length are weak predictors of detector-labeled contradiction. The maximum absolute Pearson correlation observed between the detector label and prompt length across runs is below 0.22, and between the detector label and response length below 0.20. This supports the conclusion that shallow surface features are insufficient for contradiction prediction, motivating semantic checks and contradiction heuristics.


\section*{Qualitative Error Analysis}
The qualitative review examined hallucination cases to identify common patterns across both compact reference-based prompts and the enhanced non-atomic stress test.

\begin{itemize}[leftmargin=1.5em]
\item \textbf{Confident misremembering.} Earlier endpoints, especially gpt-3.5-turbo-0125, often produced plausible and confident answers that were factually wrong. For example, when asked about a specialized term, the model sometimes gave a real term from a related domain rather than the reference answer. These cases are important because they resemble ordinary hallucinations: the output is fluent and specific, but the content conflicts with the reference.

\item \textbf{Negation and unavailability errors.} Some hallucinations involved incorrectly stating that something did not exist, was unavailable, or could not be identified when the reference contained a concrete answer. These cases are well matched to the detector's contradiction-cue and unavailability heuristics because the model response directly negates the premise or availability of the reference answer.

\item \textbf{Over-answering.} Some responses began with a correct or nearly correct core answer but then added extra details that introduced inaccuracies. Under a strict contradiction policy, a response can be flagged if an added claim conflicts with the reference or source constraint, even when the first part of the answer is useful. This remains a difficult edge case because the core answer and the added explanation may deserve separate judgments.

\item \textbf{Riddle guessing.} Many incorrect riddle responses were direct guesses with little semantic relation to the expected solution. These cases often behave like atomic-answer contradictions because the response presents a mutually exclusive answer to a prompt with a short expected solution.

\item \textbf{Computation and numeric mistakes.} A smaller set of errors involved numeric or simple reasoning mistakes. The detector is relatively sensitive to these because numeric conflict is an explicit signal. However, the audit also shows why numeric normalization matters: equivalent notation, missing units, or underspecified direction can otherwise create false positives.

\item \textbf{Abstention and calibration.} The gpt-5-2025-08-07 endpoint produced many more empty or abstaining first-line answers than the other evaluated endpoints. This behavior lowers detector-labeled contradiction-style hallucination rates because the detector treats abstention as non-contradiction. The result should therefore be interpreted as partly reflecting calibration and answer-withholding behavior, not only improved task success.

\item \textbf{Semantic inadequacy without explicit contradiction.} The enhanced non-atomic audit makes this failure mode more visible. Some answers are vague, incomplete, unsupported, or insufficiently responsive without directly contradicting the reference. These cases are not counted as positives in the binary explicit-contradiction metrics, but they are reported separately through labels such as \code{incomplete\_or\_vacuous} and \code{unsupported\_extra\_claim}.
\end{itemize}

Overall, these patterns show that hallucination is not monolithic. The benchmark targets the explicit-contradiction portion of hallucination behavior, while the expanded audit separately records semantic-adequacy failures that fall outside that binary target. This distinction is necessary for interpreting the detector's precision, recall, and F1 scores.

\section*{Discussion}
The results show lower detector-labeled contradiction rates for newer GPT endpoints on the legacy short-answer benchmark, but the enhanced audit shows that this pattern does not transfer cleanly to all non-atomic hallucination-evaluation settings. The central finding is two sided, in that explicit contradiction detection is useful under compact reference-based conditions, and its precision can degrade predictably when uncertainty expression and source limitation become part of the task.

\textbf{Conservatism: contradiction detection vs. correctness.} Because HalluDetector labels hallucination only on contradiction evidence, the metric is best interpreted as a lower bound on ``being wrong.'' A model can be wrong yet not contradict the reference strongly enough to trigger a veto, especially when the response is vague or content-free. Thus, the measured quantity is closer to ``rate of contradicted factual claims'' than ``task error rate.'' This is a reasonable definition in some safety settings, but it must be stated explicitly when comparing to accuracy-focused benchmarks.

\textbf{Re-interpreting 0\% on riddles for gpt-5-2025-08-07.} A 0--2\% detector-labeled contradiction rate for riddles under this policy does not guarantee perfect riddle-solving; it indicates that the detector did not identify contradictions. For riddles, stricter scoring would yield a more task-accuracy-like picture.

\textbf{Marginal utility of external contradiction detectors.} The high abstention rate observed for gpt-5-2025-08-07 suggests that part of the measured decline in detector-labeled contradiction reflects stronger intrinsic calibration rather than only fewer underlying reasoning errors. As frontier models become more willing to abstain or express uncertainty, the marginal value of an external contradiction detector may decrease for simple atomic prompts. However, such detectors can still be useful as an auditable safety layer for regression testing, cross-model comparisons, retrieval-grounded verification, and cases where a model gives a confident answer instead of abstaining. The results should therefore be read as evidence that the detector's role may shift over time, not as evidence that external verification becomes unnecessary.

\textbf{Practical improvements.} The enhanced-audit failures are not evidence that contradiction detection is unusable. They identify specific detector rules that can be refined. The most direct improvement is an uncertainty-aware guardrail that distinguishes asserting a false claim from mentioning uncertainty, missing evidence, unavailable sources, or claims that should be avoided. A second improvement is an assertion-scope check that separates what the model endorses from what it merely names as unsupported or prohibited. For the legacy short-answer setting, smaller refinements such as scientific-notation normalization, alias dictionaries for common-name or spelling variants, and vacuous-answer detection for atomic definitional questions remain useful. These changes can be implemented without changing the overall architecture and tested against the existing audit set.

\textbf{Layered evaluator uncertainty.} HalluDetector does not eliminate uncertainty by moving evaluation from a generative model to a detector. The detector also relies on imperfect components, including RoBERTa-MNLI, embedding similarity, lexical heuristics, numeric thresholds, and adjudication. This creates a layered evaluation setting in which the base model may contradict the reference, the detector may misclassify the response, and annotators may disagree about borderline cases. The justification for this architecture is therefore not that the detector is infallible, but that its rules, thresholds, intermediate signals, and audit disagreements are explicit and reproducible. In this study, the external detector improves reliability by converting opaque model behavior into auditable contradiction-evidence decisions and by exposing where those decisions fail, rather than by removing evaluator uncertainty entirely.

\textbf{Assumptions.} The analysis rests on several explicit assumptions. First, the reference answers are treated as sufficiently correct for reference-based contradiction scoring; this is reasonable for a controlled prompt-reference benchmark, but any reference error would propagate into detector and audit labels. Second, archived endpoint identifiers and recorded outputs are treated as the evaluated behavior during the collection window. Third, the contradiction-evidence policy is assumed to be suitable for measuring explicit contradiction risk, not general answer quality. Fourth, adjudicated audit labels are treated as final labels while agreement statistics are reported to show residual subjectivity. Fifth, repeated prompts and repeated prompt-response pairs are not assumed to be independent generalization evidence, which is why row-level metrics are reported alongside deduplicated prompt-response and prompt-identity statistics.

\textbf{Threats to validity.} These results should be interpreted in light of several scope conditions. First, the updated corpus contains 250 prompt-reference entries: the legacy 150-prompt benchmark spanning factual questions, riddles, and auto-generated adversarial prompts, plus a 100-prompt enhanced non-atomic benchmark. This remains a compact but deliberately varied corpus rather than a comprehensive estimate of general-purpose hallucination behavior. Second, the archived outputs preserve API model identifiers, decoding defaults, prompts, model responses, detector labels, and run timestamps where available; accordingly, the cross-model comparison should be read as a comparison of recorded API behavior during the collection window. Third, the auto-generated references were archived as final reference strings from the generation scripts, which supports reproducibility, though a separately logged external verification pass would strengthen future benchmark releases. Fourth, the revised 850-row audit is intentionally stratified and includes both the legacy short-answer audit rows and an enhanced non-atomic audit sample; because repeated prompts and repeated prompt-response pairs are present, row-level precision/recall values are reported together with deduplicated prompt-response metrics, prompt-identity aggregation, and prompt-clustered confidence intervals rather than treated as independent population estimates. Finally, because the detector intentionally treats incompleteness and abstention as non-hallucination unless contradiction evidence is present, the reported rates are best read as conservative contradiction-focused hallucination rates rather than a full task-error metric.

\textbf{Beyond reference-based detection.} Many deployments lack a gold reference. One path is retrieval-augmented verification: retrieve evidence from trusted corpora and treat that evidence as a proxy reference, then apply a similar contradiction detector (connected to FEVER-like pipelines) (4). Another path is model self-consistency (e.g., SelfCheckGPT) (6), trading compute for reference-free detection. Related corpus-building efforts for trustworthy hallucination evaluation include RAGTruth (12).

\section*{Conclusion}
Using the updated 250-entry prompt corpus across legacy and enhanced phases (5,000 archived model answers), the recorded outputs show lower detector-labeled contradiction rates on the legacy 150-prompt benchmark from gpt-3.5-turbo-0125 to gpt-4-0613 to gpt-4o-2024-08-06 to gpt-5-2025-08-07. An expanded 850-row adjudicated audit shows that detector performance depends strongly on prompt type: the detector remains strong on the legacy short-answer audit subset, but the enhanced non-atomic subset reveals a substantial precision drop. The detector catches most enhanced explicit contradictions, but it over-flags calibrated answers that discuss missing evidence, source limitations, claims to avoid, or cautious non-answers. This is the main boundary identified by the study. Explicit contradiction detection is not a general hallucination or semantic-adequacy detector; it is most defensible as a recall-oriented screen for compact reference disagreements. The enhanced result is informative because it shows why contradiction-based hallucination detection is entangled with uncertainty expression, abstention, partial correctness, and explanation structure. Future work should combine contradiction evidence with adequacy constraints, uncertainty-aware guardrails, and larger category-balanced non-atomic audits.

\section*{Audit Set Construction}
The audit table used for this revision contains 850 adjudicated rows. It combines 450 rows from the legacy short-answer audit with 400 rows from the enhanced audit sample. The enhanced sample was included to stress the detector with prompts that are less favorable to short-answer contradiction scoring, including longer-form explanation, causal reasoning, synthesis, citation-grounded answering, and uncertainty-oriented prompts. The enhanced prompt database contains 100 prompts, organized as 20 prompts in each of the five enhanced template categories.

Each audit row contains the source file, template, prompt ID, prompt, reference answer, model response, detector label, two annotator labels, annotator rationales, an agreement flag, adjudication information, duplicate identifiers, and prompt-cluster identifiers. The final adjudicated label is one of \code{supported}, \code{contradiction}, \code{incomplete\_or\_vacuous}, \code{abstention}, \code{unsupported\_extra\_claim}, or \code{needs\_review}. The binary explicit-contradiction variable is derived from the adjudicated label and is used for precision, recall, and F1 calculations.

Two annotators labeled the audit rows independently using the written codebook and were blinded to detector output and model identity. Disagreements were resolved by adjudication. The annotators agreed on 811 of 850 rows (95.4\%), with Cohen's $\kappa = 0.845$. The final adjudicated label distribution was: 717 \code{supported}, 69 \code{contradiction}, 57 \code{incomplete\_or\_vacuous}, 3 \code{abstention}, 3 \code{unsupported\_extra\_claim}, and 1 \code{needs\_review}.

Because repeated prompts and repeated prompt-response pairs occur in the audit, the file includes \code{prompt\_hash}, \code{response\_hash}, \code{prompt\_response\_hash}, \code{unique\_prompt\_id}, duplicate flags, prompt occurrence counts, and cluster identifiers. The adjudicated audit contains 242 underlying prompt identities represented in the sampled audit and 761 unique prompt-response pairs. Accordingly, the manuscript reports row-level metrics, deduplicated prompt-response metrics, and prompt-identity aggregation. The enhanced false-positive taxonomy was coded from the prompt, reference answer, model response, detector label, and adjudicated label.

\section*{Detector Thresholds and Defaults}
The following table lists key detector threshold values.
\begin{table}[H]
\centering
\small
\caption{Key detector thresholds.}
\begin{tabular}{p{0.28\linewidth}p{0.18\linewidth}p{0.44\linewidth}}
\toprule
Parameter & Value & Meaning\\
\midrule
numeric rel tol & 0.005 & Relative tolerance for numeric match (non-integer).\\
numeric abs tol & 0.05 & Absolute tolerance for numeric match (non-integer).\\
embed model name & sentence-transformers/all-MiniLM-L6-v2 & Sentence embedding model for cosine similarity.\\
embed min similarity & 0.72 & Cosine threshold for embedding baseline match.\\
embed min shared content tokens & 1 & Embedding gate: minimum shared content tokens.\\
jaccard min & 0.4 & Lexical Jaccard threshold when embeddings not used.\\
nli model name & roberta-large-mnli & NLI model for bidirectional contradiction.\\
nli min tokens & 3 & Minimum token length for NLI checks.\\
nli contradiction veto & 0.75 & Contradiction probability threshold for veto.\\
nli entailment veto floor & 0.3 & Entailment floor for contradiction veto.\\
atomic gold max tokens & 2 & Atomic gold heuristic: max tokens.\\
atomic resp max tokens & 5 & Atomic mismatch heuristic: max response tokens.\\
\bottomrule
\end{tabular}
\end{table}

\section*{Repo Artifacts and Reproduction Notes}
The paper references prompt CSVs, labeled outputs, run-level rates, enhanced benchmark outputs, and plots in the repository output directories. If outputs are regenerated with different endpoints or decoding settings, the tables and figures should be updated accordingly.

\section*{Acknowledgments}
The author acknowledges limited use of language-editing tools for formatting, grammar, and clarity checks during manuscript revision. The author reviewed the manuscript, audit interpretations, tables, and conclusions and is responsible for the final content.

\section*{Data and Code Availability}
The code, prompts, benchmark files, and evaluation outputs used in this study are available at \url{https://doi.org/10.6084/m9.figshare.32095732}.

\section*{References}
\small
\begin{enumerate}[leftmargin=1.5em]
\item OpenAI. GPT-4 Technical Report. arXiv, 2023. DOI: \url{https://doi.org/10.48550/arXiv.2303.08774}
\item Ouyang L, Wu J, Jiang X, Almeida D, Wainwright CL, Mishkin P, et al. Training Language Models to Follow Instructions with Human Feedback. arXiv, 2022. DOI: \url{https://doi.org/10.48550/arXiv.2203.02155}
\item Lin S, Hilton J, Evans O. TruthfulQA: Measuring How Models Mimic Human Falsehoods. ACL, 2022. DOI: \url{https://doi.org/10.18653/v1/2022.acl-long.229}
\item Thorne J, Vlachos A, Christodoulopoulos C, Mittal A. FEVER: a Large-scale Dataset for Fact Extraction and VERification. NAACL, 2018. DOI: \url{https://doi.org/10.18653/v1/N18-1074}
\item Kryscinski W, McCann B, Xiong C, Socher R. Evaluating the Factual Consistency of Abstractive Text Summarization. EMNLP, 2020. DOI: \url{https://doi.org/10.18653/v1/2020.emnlp-main.750}
\item Manakul P, Liusie A, Gales MJF. SelfCheckGPT: Zero-Resource Black-Box Hallucination Detection for Generative Language Models. EMNLP, 2023. DOI: \url{https://doi.org/10.18653/v1/2023.emnlp-main.557}
\item Reimers N, Gurevych I. Sentence-BERT: Sentence Embeddings using Siamese BERT-Networks. EMNLP-IJCNLP, 2019. DOI: \url{https://doi.org/10.18653/v1/D19-1410}
\item Liu Y, Ott M, Goyal N, Du J, Joshi M, Chen D, et al. RoBERTa: A Robustly Optimized BERT Pretraining Approach. arXiv, 2019. DOI: \url{https://doi.org/10.48550/arXiv.1907.11692}
\item Li J, Chen J, Ren R, Cheng X, Zhao X, Nie JY, et al. The Dawn After the Dark: An Empirical Study on Factuality Hallucination in Large Language Models. ACL, 2024. DOI: \url{https://doi.org/10.18653/v1/2024.acl-long.586}
\item Zhang Y, Li Y, Cui L, Cai D, Liu L, Fu T, et al. Siren's Song in the AI Ocean: A Survey on Hallucination in Large Language Models. Computational Linguistics, 2025. DOI: \url{https://doi.org/10.1162/coli.a.16}
\item Wang Y, Wang M, Manzoor MA, Liu F, Georgiev GN, Das RJ, et al. Factuality of Large Language Models: A Survey. EMNLP, 2024. DOI: \url{https://doi.org/10.18653/v1/2024.emnlp-main.1088}
\item Niu C, Wu Y, Zhu J, Xu S, Shum K, Zhong R, et al. RAGTruth: A Hallucination Corpus for Developing Trustworthy Retrieval-Augmented Language Models. arXiv, 2024. DOI: \url{https://doi.org/10.48550/arXiv.2401.00396}
\end{enumerate}

\end{document}
